\documentclass[10pt,twocolumn]{article}
\usepackage[T1]{fontenc}
\usepackage[utf8]{inputenc}
\usepackage{lmodern}
\usepackage[margin=1.8cm,top=2.4cm,bottom=2cm,columnsep=0.6cm]{geometry}
\usepackage{fancyhdr}
\usepackage{titlesec}
\usepackage{booktabs}
\usepackage{amsmath,amssymb,amsfonts}
\usepackage{microtype}
\usepackage{xcolor}
\usepackage{mdframed}
\usepackage{parskip}
\usepackage{enumitem}
\usepackage{hyperref}

\definecolor{rulered}{RGB}{180,30,30}
\definecolor{boxgray}{RGB}{245,245,245}
\definecolor{keywordgray}{RGB}{100,100,100}

\pagestyle{fancy}
\fancyhf{}
\fancyhead[L]{\textsc{\small Claude Mag}}
\fancyhead[C]{\small\textit{Machine Learning Research Digest}}
\fancyhead[R]{\small 2026-10-28}
\fancyfoot[C]{\small\thepage}
\renewcommand{\headrulewidth}{0.8pt}

\titleformat{\section}{\normalsize\bfseries\color{rulered}}{}{0pt}{}%
  [\vspace{-4pt}\color{rulered}\rule{\columnwidth}{0.4pt}]
\titlespacing{\section}{0pt}{8pt}{4pt}

\hypersetup{colorlinks=true,urlcolor=rulered,linkcolor=black}

\begin{document}
\setlength{\parindent}{0pt}
\setlength{\parskip}{4pt}

{\small\color{keywordgray}\textbf{Keywords:}
  video diffusion \textbullet{}
  novel view synthesis \textbullet{}
  autonomous driving simulation \textbullet{}
  geometry grounding}\par\vspace{4pt}

{\LARGE\bfseries\raggedright
  GeoV2V: Geometry-Grounded Video Diffusion Model
  for Driving Scene Generation}\par\vspace{4pt}

{\large\itshape\raggedright
  LiDAR-Warped Reference Frames and Synthetic Para4D Data Enable
  Zero-Shot Lane-Shift Video Generation on Real-World Waymo
  and nuScenes Benchmarks}\par\vspace{6pt}

{\small\textbf{By} Yuchen Xi, Yuhan Guo, Chenwei Hou, Zixu Liu,
  Jin Fang, Jason Liu, Ruigang Yang}\par\vspace{2pt}

{\small\color{keywordgray}arXiv:2610.25245 \textbullet{}
  ECCV 2026}
\par\vspace{2pt}\noindent{\color{rulered}\rule{\columnwidth}{0.6pt}}\vspace{6pt}

Autonomous driving simulators require realistic video from camera
poses that were never recorded---a lane shift, a slightly different
heading, or a hypothetical trailing vehicle. GeoV2V addresses this
challenge by combining LiDAR and depth-based geometric priors with
full-video conditioning on the Wan~2.1 video diffusion backbone,
generating temporally consistent, geometrically accurate driving
videos under large trajectory shifts. Trained entirely on Para4D,
a new synthetic CARLA-based dataset of synchronized multi-lane videos,
GeoV2V transfers zero-shot to real-world Waymo and nuScenes data,
achieving state-of-the-art visual quality and geometric consistency.

\begin{mdframed}[backgroundcolor=boxgray,linecolor=rulered,linewidth=1pt,
    innerleftmargin=6pt,innerrightmargin=6pt,
    innertopmargin=6pt,innerbottommargin=6pt]
\textbf{\small AT A GLANCE}\par\vspace{4pt}
\begin{description}[leftmargin=0pt,labelsep=4pt,itemsep=2pt,parsep=0pt,
    font=\small\bfseries]
  \item[Problem:] Large viewpoint shifts in novel-view driving video
    reveal occluded regions; pure generation loses geometric coherence
    across frames.
  \item[Why it matters:] High-fidelity novel-view synthesis is essential
    for scalable simulation of rare events (lane changes, near-misses)
    that cannot be captured by fixed sensor rigs.
  \item[Method:] Geometric warping of reference frames via LiDAR/depth
    priors + Wan 2.1 video diffusion conditioned on warped sequence +
    Para4D synthetic training data.
  \item[Result:] State-of-the-art FID, FVD, and depth-consistency on
    Waymo and nuScenes without fine-tuning on real data.
  \item[Evidence:] Waymo, nuScenes, and Para4D benchmarks; ablations
    isolating geometric conditioning and synthetic pretraining.
\end{description}
\end{mdframed}

\section{The Big Picture}

The long tail of driving scenarios---rare lane changes, unexpected
obstacles, adverse weather---cannot be covered by fixed-route data
collection. Novel view synthesis (NVS) addresses this by generating
plausible videos from camera poses never recorded. For autonomous
driving, NVS must satisfy two requirements that pure appearance-based
generation fails to meet: (1) geometric consistency across time, and
(2) faithfulness to 3D scene structure rather than statistical plausibility.

GeoV2V grounds a powerful video diffusion prior (Wan~2.1) with
explicit 3D geometry, turning the generation problem into a
geometry-aware inpainting problem: warp what is visible, inpaint
what is occluded, and preserve temporal dynamics.

\section{Why Existing Approaches Fall Short}

\textbf{Image-based NVS methods (NeRF, 3DGS):}
These methods require dense multi-view captures or long sequences for
3D reconstruction---unavailable in drive-by single-camera settings.

\textbf{Pure video generation (Sora, Gen-3):}
Unconstrained video generation ignores 3D structure, producing
temporally inconsistent lane layouts and incorrect depth ordering.

\textbf{Warping-only methods:}
Geometric warp alone leaves holes where occluded content is revealed
by the viewpoint shift; filling them with inpainting lacks temporal
coherence across the full video.

\section{Core Mechanism}

GeoV2V operates in three stages per generated video:

\begin{enumerate}[itemsep=2pt,parsep=0pt]
  \item \textbf{Geometric warping:} Given the reference video and a
    target viewpoint trajectory, LiDAR point clouds (or monocular depth
    where LiDAR is unavailable) project each reference frame to the
    target pose, producing a warped reference sequence with explicit
    holes where geometry is occluded.
  \item \textbf{Conditional video generation:} The Wan~2.1 video
    diffusion transformer takes the warped reference sequence as a
    conditioning signal via cross-attention, filling holes and refining
    geometric artifacts while preserving large-scale 3D structure.
  \item \textbf{Para4D training:} The model is trained on Para4D, a
    synthetic multi-lane CARLA dataset providing pairs of synchronized
    cameras at different lateral offsets under diverse weather, lighting,
    and traffic conditions---the only source of explicit cross-lane
    supervision.
\end{enumerate}

\section{Technical Formulation}

Let $V_\mathrm{ref} = \{I_t^{\mathrm{ref}}\}_{t=1}^T$ be the reference
video and $P_\mathrm{tgt}$ the target camera pose sequence. The warp
operator $\mathcal{W}$ uses depth $D_t$ and camera intrinsics $K$ to
project pixels:
\[
\tilde{I}_t^{\mathrm{tgt}} = \mathcal{W}(I_t^{\mathrm{ref}}, D_t,
  P_t^{\mathrm{ref}} \to P_t^{\mathrm{tgt}}, K)
\]
yielding a warped sequence $\tilde{V} = \{\tilde{I}_t^{\mathrm{tgt}}\}$
with per-pixel validity masks $M_t$ marking visible regions.

The video diffusion model $G_\theta$ generates the output:
\[
V_\mathrm{out} = G_\theta\!\left(z_T, \tilde{V}, M\right)
\]
where $z_T$ is Gaussian noise and conditioning on $(\tilde{V}, M)$
is implemented via cross-attention at each transformer layer.
The training objective is the standard flow-matching loss
$\mathcal{L} = \mathbb{E}[\|G_\theta(z_t, \tilde{V}, M) - V_\mathrm{gt}\|^2]$
on Para4D pairs.

\section{Learning or Inference Procedure}

\textbf{Training:} GeoV2V is trained entirely on synthetic Para4D data
(no real-world data). The model sees diverse CARLA scenes with
lateral offsets of 1.0--5.0\,m, covering standard lane shifts and
extreme viewpoints. Wan~2.1 weights are initialised from a pretrained
checkpoint and fine-tuned with the geometric conditioning added.

\textbf{Inference (zero-shot transfer):} On Waymo and nuScenes, only
LiDAR-based depth is used for warping; the generator runs without
any domain-specific fine-tuning, relying on the synthetic-to-real
generalisation of the learned geometric prior.

\section{What the Guarantee Says}

The geometric warping step provides an explicit constraint: visible
pixels in the reference frame are reproduced at the correct target
location with sub-pixel registration accuracy (verified by LiDAR
reprojection error). The diffusion model only hallucrinates content
in provably occluded regions (mask $M_t = 0$), giving a formal
separation between geometric reliability and generative uncertainty.

\section{Experimental Findings}

\begin{itemize}[itemsep=2pt,parsep=0pt]
  \item \textbf{Waymo lane-shift NVS:} State-of-the-art FID and FVD
    against prior methods (WoVoGen, GeoFlow, AnyScene), with
    $>$15\% improvement in FVD at 3.5\,m lane shifts.
  \item \textbf{nuScenes:} Consistent gains in geometric depth
    consistency (homography error reduction $>$20\%) and visual
    quality metrics.
  \item \textbf{Challenging lighting:} GeoV2V maintains FVD gains in
    night and adverse weather conditions where geometric priors are
    noisiest.
  \item \textbf{Para4D ablation:} Replacing Para4D with a warping-only
    baseline increases FVD by 38\%, confirming the importance of
    explicit cross-lane supervision.
\end{itemize}

\section{Ablations and Interpretation}

\textbf{LiDAR vs.\ monocular depth:}
LiDAR warping improves depth consistency by 12\% over monocular
depth; monocular depth is sufficient for visual quality (FVD) but
not for geometric accuracy (reprojection error).

\textbf{Conditioning mechanism:}
Replacing cross-attention conditioning with concatenation reduces FVD
by 8\%, confirming that the temporal context of the warped sequence
is best exploited through attention.

\textbf{Training data scale:}
Para4D with 10k scene pairs is sufficient; doubling to 20k improves
FVD by only 2\%, suggesting diminishing returns from synthetic data.

\section*{Reference}

Yuchen Xi, Yuhan Guo, Chenwei Hou, Zixu Liu, Jin Fang, Jason Liu,
Ruigang Yang.
\textbf{GeoV2V: Geometry-Grounded Video Diffusion Model for Driving
Scene Generation}.
arXiv:2610.25245. ECCV 2026.
\url{https://arxiv.org/abs/2610.25245}

\end{document}
